% Baccalauréat malgache, série C, 2005 : énoncé.
\TCDBacAttribution{C}{2005}
\providecommand{\dx}{\mathrm{d}x}
\providecommand{\Re}[1]{\mathrm{Re}\left( #1 \right)}
\providecommand{\Im}[1]{\mathrm{Im}\left( #1 \right)}
\providecommand{\dt}{\mathrm{d}t}
\providecommand{\Ouv}{(O;\vec{u},\vec{v})}
\providecommand{\mes}[2]{\mathrm{mes}\left( \widehat{ \overrightarrow{#1} ; \overrightarrow{#2} } \right)}
\providecommand{\D}{(\mathscr{D})}
\providecommand{\Courbe}{(\mathscr{C})}
\providecommand{\Courben}{(\mathscr{C}_n)}
\providecommand{\CourbeU}{(\mathscr{C}_1)}
\providecommand{\CourbeD}{(\mathscr{C}_2)}
\providecommand{\CourbeA}{(\mathscr{C}_\alpha)}
\providecommand{\congru}[3]{#1 \equi #2 \quad \left[ \mathrm{mod}~#3\right]}
\providecommand{\Angl}[2]{\left( \widehat{ \overrightarrow{#1} , \overrightarrow{#2} } \right)}
%
%%%%%%%%%%%%%%%%%%%%%%%%%%%%%%%%%%%%%%%%%%%%%%%%%%%%%%%%%%%%%%%%%%%%%%%%%%%%%%%%%%%%%%%%%%%%%%%%%%%%%%%%%%%%%%%%%%%%%%%%%%%%%%%%
\noindent \underline{\textsc{\textbf{Exercice}}} (4 points) 
      \begin{enumerate}
            \item On donne l'\'equation : 
                  \begin{equation}
                        \TCDtag{1}
                              2 x - y = 1 \; ;
                        \label{eq:Bac C 2005 Arithm 1}
                  \end{equation}
                  $x$ et $y$ \'etant les inconnues, $x$ et $y$ sont des entiers relatifs.
                  \begin{enumerate}
                        \item Trouver une solution particuli\`ere $(x_0 ; y_0)$ de cette \'equation.
                        \item R\'esoudre dans $\mathbb{Z} \times \mathbb{Z}$ l'\'equation \eqref{eq:Bac C 2005 Arithm 1}.
                  \end{enumerate}
            \item Un entier naturel $A$ est tel que :
                  \begin{itemize}
                        \item dans la base 10, $A$ s'\'ecrit $A=b+6$ ;
                        \item dans le syst\`eme binaire, $A$ s'\'ecrit $A=\left(\overline{1a1}\right)_2$.
                  \end{itemize}
                  
                  Trouver une relation qui lie $a$ et $b$.
            \item R\'esoudre dans $\mathbb{N}^2$ le syst\`eme suivant, d'inconnues $x$ et $y$ :
                  \[
                        \begin{multiform}
                              x+y=60 \\ \pgcd{x}{y}=12.
                        \end{multiform}
                  \]
            \item Deux urnes identiques $U_1$ et $U_2$ contiennent chacune : 
                  \begin{itemize}
                        \item deux boules num\'erot\'ees chacune 0 ;
                        \item trois boules num\'erot\'ees chacune 1 ;
                        \item une boule num\'erot\'ee 2.
                  \end{itemize}
                  
                  On tire au hasard une boule de $U_1$, puis on extrait au hasard et simultan\'ement deux boules de $U_2$.
                  \begin{enumerate}
                        \item Calculer la probabilit\'e de chacun des \'ev\`enements suivants : 
                              \begin{itemize}
                                    \item $A$ : << Les boules obtenues portent le m\^eme num\'ero >>.
                                    \item $B$ : << La somme des num\'eros not\'es est \'egale \`a 4 >>.
                              \end{itemize}
                        \item Soient $x$ le num\'ero de la premi\`ere boule tir\'ee et $y$ la somme des num\'eros des 
                              boules obtenues au deuxi\`eme tirage.
                              Calculer la probabilit\'e de l'\'ev\`enement : 
                              $C$ : << Le couple $(x,y)$ est solution de l'\'equation $2 x - y = 1$ >>.
                  \end{enumerate}
      \end{enumerate}
%%%%%%%%%%%%%%%%%%%%%%%%%%%%%%%%%%%%%%%%%%%%%%%%%%%%%%%%%%%%%%%%%%%%%%%%%%%%%%%%%%%%%%%%%%%%%%%%%%%%%%%%%%%%%%%%%%%%%%%%%%%%%%%%
\noindent \underline{\textsc{\textbf{PROBLEME I}}} (7 points) \newline
      Soit $(ABC)$ un triangle isoc\`ele et rectangle en $A$ avec $AB = AC = 4\cm$ et $\Angl{AB}{AC} = \frac{\pi}{2}$.
      Soient $I$, $J$ et $K$ les milieux respectifs des segments $[AB]$, $[BC]$ et $[CA]$ et $E$ le  sym\'etrique de $A$ 
      par rapport \`a $J$.
      \begin{description}
            \item[Partie 1]
                  \begin{enumerate}
                        \item \begin{enumerate}
                                    \item Montrer que $\overrightarrow{JI} = \frac{\overrightarrow{JA} + \overrightarrow{JB}}{2}$
                                          et $\overrightarrow{JK} = \frac{\overrightarrow{JA} + \overrightarrow{JC}}{2}$.
                                          En d\'eduire que $\overrightarrow{IK} = \frac{1}{2} \overrightarrow{BC}$.
                                    \item Montrer que $\overrightarrow{KI} = \frac{\overrightarrow{KA} + \overrightarrow{KB}}{2}$
                                          et $\overrightarrow{KJ} = \frac{\overrightarrow{KB}+\overrightarrow{KC}}{2}$.
                                          En d\'eduire que $\overrightarrow{IJ} = \frac{1}{2} \overrightarrow{AC}$.
                                    \item Montrer que le quadruplet $(AIJK)$ est un carr\'e. Faire une figure.
                              \end{enumerate}
                        \item Soit $r$ la rotation de centre $J$ et d'angle $\frac{\pi}{2}$ et $t$ la translation de 
                              vecteur $\vec{u} = \frac{1}{2} \overrightarrow{BC}$.
                              \begin{enumerate}
                                    \item D\'eterminer l'image du point $I$ par $t$ et celle de $K$ par $r$.
                                    \item D\'eterminer l'image de $I$ par $f = r \circ t$. 
                                          En d\'eduire la nature et les \'el\'ements caract\'eristiques de $f$.
                                    \item Soit $g = t \circ r$.
                                          D\'eterminer l'image de $K$ par $g$.
                                          En d\'eduire la nature et les \'el\'ements caract\'eristiques de $g$.
                              \end{enumerate}
                        \item Soit $S$ la similitude plane directe qui transforme $B$ en $E$ et $J$ en $C$.
                              Pr\'eciser le rapport et l'angle de $S$.
                              \label{quest: Bac C 2005 Pbm 1 Part 1 quest 3}
                              
                              \underline{N.B :} Justifier votre r\'eponse.
                  \end{enumerate}
            \item[Partie 2]
                  Le plan complexe est rapport\'e \`a un rep\`ere orthonorm\'e $(A ; \vec{e_1} , \vec{e_2})$ o\`u 
                  $\vec{e_1} = \frac{1}{4} \overrightarrow{AB}$ et $\vec{e_2} = \frac{1}{4} \overrightarrow{AC}$.
                  Dans ce rep\`ere, on donne les points : $A(0,0)$, $B(4,0)$ et $C(0,4)$.
                  \begin{enumerate}
                        \item Pr\'eciser les coordonn\'ees des points $E$ et $J$.
                        \item \begin{enumerate}
                                    \item Donner l'\'ecriture complexe de $S$.
                                          ($S$ \'etant la similitude plane directe donn\'ee dans la \textbf{Partie 1},
                                          \ref{quest: Bac C 2005 Pbm 1 Part 1 quest 3})
                                    \item En d\'eduire les coordonn\'ees du centre $\Omega$ de $S$.
                              \end{enumerate}
                        \item Soit la transformation $\overline{S} = S \circ S_{(AB)}$, o\`u $S_{(AB)}$ est la sym\'etrie
                              orthogonale par rapport \`a la droite $(AB)$.
                              \begin{enumerate}
                                    \item Pr\'eciser la nature et les \'el\'ements caract\'eristiques de $\bar{S}$.
                                    \item D\'eterminer l'\'ecriture complexe associ\'ee \`a $\bar{S}$.
                              \end{enumerate}
                  \end{enumerate}
      \end{description}
%%%%%%%%%%%%%%%%%%%%%%%%%%%%%%%%%%%%%%%%%%%%%%%%%%%%%%%%%%%%%%%%%%%%%%%%%%%%%%%%%%%%%%%%%%%%%%%%%%%%%%%%%%%%%%%%%%%%%%%%%%%%%%%%
\noindent \underline{\textsc{\textbf{PROBLEME II}}} (9 points) \newline
      Soit la fonction num\'erique $f$ d\'efinie sur $\mathbb{R}$ par : 
      \[
            f(x) = 
            \begin{multiform}
                  (1-x) \expo{\frac{x}{2}} \quad \text{ si } x \geq 0 \\
                  1 + \frac{x}{2} \ln (1-x) \quad \text{ si } x < 0 \;.
            \end{multiform}
      \]
      $\Courbe$ sa repr\'esentation graphique dans un plan affine rapport\'e \`a un rep\`ere orthonorm\'e.
      \begin{description}
            \item[Partie A]
                  \begin{enumerate}
                        \item Soit $g$ la fonction d\'efinie sur l'intervalle $]-\infty ; 0[$ par : 
                              $g(x) = (1-x) \ln (1-x) - x$.
                              \begin{enumerate}
                                    \item Etudier les variations de $g$ et dresser son tableau de variation.
                                    \item En d\'eduire le signe de $g(x)$ suivant les valeurs de $x$.
                              \end{enumerate}
                        \item \begin{enumerate}
                                    \item Montrer que $f$ est continue en $x_0 = 0$.
                                    \item Calculer $\displaystyle{ \lim_{x\rightarrow 0^-} \frac{f(x) - f(0)}{x} }$.
                                          \label{quest: Bac C 2005 Pbm 2 Part A 2/ b}
                                    \item Pour tout $x > 0$, montrer que
                                          \[
                                                \frac{f(x) - f(0)}{x} = 
                                                \frac{1}{2}
                                                \left[ \frac{ \expo{\frac{x}{2}}-1 }{ \frac{x}{2}} \right] -
                                                \expo{\frac{x}{2}}
                                          \]
                                          Calculer 
                                          \[
                                          \lim_{x\rightarrow 0^+ } \frac{f(x) - f(0)}{x}
                                          \]
                                          (Indication : poser $X = \frac{x}{2}$.)
                                          \label{quest: Bac C 2005 Pbm 2 Part A 2/ c}
                                    \item En utilisant les r\'esultats de 
                                          \ref{quest: Bac C 2005 Pbm 2 Part A 2/ b} et \ref{quest: Bac C 2005 Pbm 2 Part A 2/ c},
                                          quelles conclusions peut-on tirer sur $f$ puis sur $\Courbe$\,?
                              \end{enumerate}
                        \item \begin{enumerate}
                                    \item Calculer 
                                          $\displaystyle{
                                          \lim_{x\rightarrow -\infty } f(x)
                                          }$
                                          et
                                          $\displaystyle{
                                          \lim_{x\rightarrow +\infty } f(x)
                                          }$.
                                    \item Pour tout $x$ \'el\'ement de l'intervalle $]-\infty ; 0[$, montrer que : 
                                          $f'(x) = \frac{g(x)}{2 (1-x)}$.
                                    \item Pour tout $x$ élément de l’intervalle $]0;+\infty[$, calculer $f'(x)$.
                                    \item Dresser le tableau de variation de $f$.
                              \end{enumerate}
                        \item Etudier les branches infinies de $\Courbe$.
                        \item Montrer que la courbe $\Courbe$ coupe l'axe des abscisses en deux points distincts 
                              d'abscisses respectives $1$ et $\alpha$ avec $-2 < \alpha < -1$.
                        \item Tracer la courbe $\Courbe$ en pr\'ecisant les demi-tangentes \`a l'origine $O$ du rep\`ere.
                        \item Calculer, en $\cm^2$, l'aire $A$ du domaine plan limit\'e par $\Courbe$, l'axe des
                              abscisses et les droites d'\'equation $x = 0$ et $x = 1$.
                  \end{enumerate}
            \item[Partie B]
                  Soit $(I_n)_{n\in\mathbb{N}}$ la suite d\'efinie par : $I_0 = \sqrt{e} - 1$ et 
                  pour tout $n \in \mathbb{N}^*$, 
                  $I_n = \frac{1}{2^{n+1} (n!)} \int_{0}^{1} (1-t)^n \; \expo{\frac{t}{2}} \dt$.
                  \begin{enumerate}
                        \item A l'aide d'une int\'egration par parties, calculer $I_1$.
                        \item Montrer que pour tout $n$ \'el\'ement de $\mathbb{N}$, on a 
                              $I_{n+1} = I_n - \frac{1}{2^{n+1} [(n+1)!]}$.
                  \end{enumerate}
            \item En d\'eduire par r\'ecurrence que, pour tout entier naturel $n$, on a
                  \[
                  \sqrt{e}=1+\frac{1}{2}\frac{1}{1!}
                            +\frac{1}{2^2}\frac{1}{2!}
                            +\cdots
                            +\frac{1}{2^n}\frac{1}{n!}+I_n.
                  \]
            \item Donner alors l'expression de $I_n$ en fonction de $n$.
      \end{description}
      \begin{center}
            On donne : $\ln 2 \simeq 0,7$, $\ln 3 \simeq 1,1$, $\sqrt{e} \simeq 1,64$.
      \end{center}
%%%%%%%%%%%%%%%%%%%%%%%%%%%%%%%%%%%%%%%%%%%%%%%%%%%%%%%%%%%%%%%%%%%%%%%%%%%%%%%%%%%%%%%%%%%%%%%%%%%%%%%%%%%%%%%%%%%%%%%%%%%%%%%%


































